% !TEX root = ../../../main.tex
% Block card: detection.fluss -- "FLUSS regime change" (Part I, the blocks chapter).
\cardfield{What it does}
Scores every moment of the span by how strongly the recording seems to change its kind of
behaviour there: 0 for no sign of a change, towards 1 for a clear one.

\begin{plainterms}[How FLUSS sees a change]
The matrix profile (\blockref{detection.matrix_profile}) finds, for every window, its closest
look-alike. Draw a string from each window to its look-alike. While the recording keeps doing the
same kind of thing, the strings stay within that stretch; where it switches to something else,
few strings cross the switch, because nothing before it looks like what comes after. Count the
strings passing over each moment, compare with how many would pass there by chance, and a dip in
that count is a change. The score is high where the dip is deep.
\end{plainterms}

\cardfield{When to use it}
When you suspect a recording moves between distinct regimes (say, quiet and active phases) and
want to find where, without describing the regimes in advance. Follow it with
\blockref{detection.threshold} for one span per change. It is marked experimental: on the real
recordings tried so far it found nothing beyond what shuffled copies of the same recordings give.

\cardfield{What goes in, what comes out}
In goes a span of one channel. Out comes one score per sample, between 0 and 1. The two ends of the
span, five windows each at the default, and the last window's worth of samples carry no score.

\cardfield{Settings that matter}
\emph{window\_s} is the window in seconds: set it to about one period of the behaviour, one typical
event or cycle, and tune it first. \emph{excl\_factor} is how many windows are left unscored at
each end. \emph{n\_regimes\_readout} only changes how many candidate changes the block lists in its
numbers; it does not change the score.

\cardfield{Watch out for}
Detrend first: on a drifting recording every window looks like every other and the score stays
flat. A change near either end of the span cannot be seen. A threshold chosen for one window
length does not hold at another, because longer windows give higher peaks everywhere. It costs a
full matrix profile, so long spans are slow.
